
\subsection{定义 $A_\infty$ 双模同态的障碍 II}
\label{subsec:5-2-7}

为完成定理 5.2.35 的证明，
现在使用能量滤过，
因而需要与第 4.5.4、4.5.5 小节
类似的论证。
障碍可借助下文定义的一种
Hochschild 上同调来表述。

令 $\overline C_i,\overline C'_i$
（$i=0,1$）为 $A_\infty$ 代数，
$\overline{\mathfrak f}^{(i)}:
\overline C_i\to\overline C'_i$
为 $A_\infty$ 同态，
并令 $\overline D,\overline D'$
为分次 $R$-模。
记
\begin{equation}
\begin{aligned}
B\overline D B
&=B\overline C_1[1]\otimes
\overline D[1]\otimes
B\overline C_0[1],\\
B\overline D'B
&=B\overline C'_1[1]\otimes
\overline D'[1]\otimes
B\overline C'_0[1].
\end{aligned}
\tag{5.2.50}\label{eq:5-2-50}
\end{equation}
按下式，
$B\overline D B$ 成为
$(B\overline C_1[1],
  B\overline C_0[1])$-双余模：
\[
\begin{aligned}
\Delta_{\mathrm{Left}}
(\mathbf x\otimes y\otimes\mathbf z)
&=\sum_a
\mathbf x_a^1\otimes
(\mathbf x_a^2\otimes y\otimes\mathbf z),
\\
\Delta_{\mathrm{Right}}
(\mathbf x\otimes y\otimes\mathbf z)
&=\sum_b
(\mathbf x\otimes y\otimes\mathbf z_b^1)
\otimes\mathbf z_b^2,
\end{aligned}
\]
其中
\begin{equation}
\Delta(\mathbf x)
=\sum_a\mathbf x_a^1\otimes\mathbf x_a^2,
\qquad
\Delta(\mathbf z)
=\sum_b\mathbf z_b^1\otimes\mathbf z_b^2.
\tag{5.2.51}\label{eq:5-2-51}
\end{equation}

\begin{definition}{5.2.52}\label{def:5-2-52}
记
\[
\begin{aligned}
&\operatorname{CH}
(\overline C_1,\overline D,\overline C_0;
 \overline C'_1,\overline D',\overline C'_0)
\\
&\quad=
\prod_{k_1=0}^{\infty}
\prod_{k_0=0}^{\infty}
\operatorname{Hom}\bigl(
B_{k_1}\overline C_1[1]\otimes
\overline D[1]\otimes
B_{k_0}\overline C_0[1],
\overline D'[1]\bigr).
\end{aligned}
\]
其中的每个元素
\[
\mathcal T
=(\mathcal T_{0,0},\ldots,
  \mathcal T_{k_1,k_0},\ldots)
\in
\operatorname{CH}
(\overline C_1,\overline D,\overline C_0;
 \overline C'_1,\overline D',\overline C'_0)
\]
都唯一延拓成双余导子\liftIndex{I124}
\[
\widehat{\mathcal T}:
B\overline D B\longrightarrow
B\overline D'B,
\]
其作用为
\translatornote{原书实体页 304
在下式及左右余作用相容式中
把左侧代数同态写为
$\widehat{\mathfrak f}$、
右侧写为
$\widehat{\mathfrak f}'$，
但本定义前只引入
$\overline{\mathfrak f}^{(1)}$、
$\overline{\mathfrak f}^{(0)}$。
这里按两侧输入的型别明确标为
$(1)$ 与 $(0)$；见疑误 A-0287。}
\[
\begin{aligned}
&\widehat{\mathcal T}
(\mathbf x,y,\mathbf z)
\\
&\quad=
\sum_{a,b,k_1,k_0}
\widehat{\overline{\mathfrak f}}^{(1)}
(\mathbf x_a^1)
\otimes
\mathcal T_{k_1,k_0}
(\mathbf x_a^2,y,\mathbf z_b^1)
\otimes
\widehat{\overline{\mathfrak f}}^{(0)}
(\mathbf z_b^2).
\end{aligned}
\]
这里使用式 (5.2.51) 的约定。
于是
\[
\begin{aligned}
\Delta_{\mathrm{Left}}\circ
\widehat{\mathcal T}
&=
\bigl(
\widehat{\overline{\mathfrak f}}^{(1)}
\otimes\widehat{\mathcal T}
\bigr)
\circ\Delta_{\mathrm{Left}},
\\
\Delta_{\mathrm{Right}}\circ
\widehat{\mathcal T}
&=
\bigl(
\widehat{\mathcal T}
\otimes
\widehat{\overline{\mathfrak f}}^{(0)}
\bigr)
\circ\Delta_{\mathrm{Right}}.
\end{aligned}
\]
若 $\overline D,\overline D'$
分别还带有所述两对代数上的
$A_\infty$ 双模结构，
则在复形
\[
\operatorname{CH}
(\overline C_1,\overline D,\overline C_0;
 \overline C'_1,\overline D',\overline C'_0)
\]
上定义上边界算子
$\delta_{\overline{\mathfrak f}^{(1)},
         \overline{\mathfrak f}^{(0)}}$
\translatornote{原书实体页 303
只把 $\overline D,\overline D'$
引入为分次 $R$-模，
实体页 304 却直接用
$\widehat d,\widehat d'$
定义边界算子。
这两个双余导子还需要相应的
$A_\infty$ 双模结构；见疑误 A-0290。}，
令
\[
\widehat{
\delta_{\overline{\mathfrak f}^{(1)},
       \overline{\mathfrak f}^{(0)}}
\mathcal T}
=
\widehat{\overline d}'\circ
\widehat{\mathcal T}
-(-1)^{\deg\mathcal T}
\widehat{\mathcal T}\circ
\widehat{\overline d}.
\]
\end{definition}

\begin{definition}{5.2.53}\label{def:5-2-53}
称复形
\[
\bigl(
\operatorname{CH}
(\overline C_1,\overline D,\overline C_0;
 \overline C'_1,\overline D',\overline C'_0),
\delta_{\overline{\mathfrak f}^{(1)},
        \overline{\mathfrak f}^{(0)}}
\bigr)
\]
为 \emph{Hochschild 复形}，
记其上同调为
\[
\operatorname{HH}
(\overline C_1,\overline D,\overline C_0;
 \overline C'_1,\overline D',\overline C'_0;
 \overline{\mathfrak f}^{(1)},
 \overline{\mathfrak f}^{(0)}),
\]
称为 \emph{Hochschild 上同调}\liftIndex{I338}。
\end{definition}

命题 4.4.48 可按显然方式推广。

\begin{definition}{5.2.54}\label{def:5-2-54}
若 $\varphi=\{\varphi_{k_1,k_0}\}$
满足定义 3.7.7 (2) 的条件，且
\translatornote{原书实体页 304
下式与上文 Hochschild 边界算子式
都未在左复合的双余导子上
区分靶端撇号。
由 $\widehat\varphi$、
$\widehat{\mathcal T}$
的源靶型别，
左项应取 $\widehat d'$，
右项取 $\widehat d$；
见疑误 A-0291。}
\[
\widehat d'\circ
\widehat\varphi_{k_1,k_0}
-
\widehat\varphi_{k_1,k_0}
\circ\widehat d
\equiv0\pmod{T^\lambda},
\]
则称它为位于
$(\mathfrak f^{(1)},\mathfrak f^{(0)})$
之上的\emph{模 $T^\lambda$
的滤过 $A_\infty$ 双模同态}\liftIndex{I323}。
\end{definition}

继续使用定义 5.2.1 的记号。
取 $G$，使这里讨论的
滤过 $A_\infty$ 代数、同态和双模
均为 $G$-有隙。
当 $G$ 无限时，可将其元素枚举为
\[
G=\{(\lambda_i,\mu_i)
\mid i=0,1,2,\ldots\},
\qquad
\lambda_0=0,\quad
\lambda_i\uparrow\infty.
\]
若 $G$ 有限，则只枚举现有层，
归纳在最后一层停止
\translatornote{原书实体页 304
把 $G$ 一律写作
$i=0,1,2,\ldots$ 的无限序列且要求
$\lambda_i\uparrow\infty$，
但定义 3.1.6 允许有限 $G$；
有限情形只需有限步。
与疑误 A-0268 同型；
见疑误 A-0289。}。
令
$\varphi=\{\varphi_{k_1,k_0}\}$
为位于
$(\mathfrak f^{(1)},\mathfrak f^{(0)})$
之上的模 $T^{\lambda_i}$
的滤过 $A_\infty$ 双模同态。
尤其，其约化
$\overline\varphi
=\{\overline\varphi_{k_1,k_0}\}$
是位于
$(\overline{\mathfrak f}^{(1)},
  \overline{\mathfrak f}^{(0)})$
之上的 $A_\infty$ 双模同态。

\begin{proposition}{5.2.55}\label{prop:5-2-55}
存在
\liftIndex{I350}
\[
o_{\lambda_{i+1}}(\varphi)
\in\operatorname{CH}
(\overline C_1,\overline D,\overline C_0;
 \overline C'_1,\overline D',\overline C'_0)
\]
使得：
\begin{enumerate}[label=\textup{(5.2.56.\arabic*)},leftmargin=5.8em]
\item\label{cond:5-2-56-1}
$\delta_{\overline{\mathfrak f}^{(1)},
          \overline{\mathfrak f}^{(0)}}
(o_{\lambda_{i+1}}(\varphi))=0$。
\item\label{cond:5-2-56-2}
$\varphi$ 可延拓为模
$T^{\lambda_{i+1}}$ 的
滤过 $A_\infty$ 双模同态，
当且仅当
\[
[o_{\lambda_{i+1}}(\varphi)]
=0\in
\operatorname{HH}
(\overline C_1,\overline D,\overline C_0;
 \overline C'_1,\overline D',\overline C'_0;
 \overline{\mathfrak f}^{(1)},
 \overline{\mathfrak f}^{(0)}).
\]
\item\label{cond:5-2-56-3}
上同调类
$[o_{\lambda_{i+1}}(\varphi)]$
只依赖于 $\varphi$ 的同伦类。
此陈述可按注 4.5.18 的方式
精确定义。
\item\label{cond:5-2-56-4}
设
\[
\overline{\mathfrak g}_{(i)}:
\overline C'_i\to\overline C''_i,
\qquad
\overline{\mathfrak g}'_{(i)}:
\overline C'''_i\to\overline C_i
\]
为 $A_\infty$ 同态，
$\overline\varphi':
\overline D'\to\overline D''$
与
$\overline\varphi'':
\overline D'''\to\overline D$
分别是位于
$(\overline{\mathfrak g}_{(1)},
  \overline{\mathfrak g}_{(0)})$
与
$(\overline{\mathfrak g}'_{(1)},
  \overline{\mathfrak g}'_{(0)})$
之上的 $A_\infty$ 双模同态。
则
\[
\bigl[
o_{\lambda_{i+1}}
(\overline\varphi'\circ
 \overline\varphi\circ
 \overline\varphi'')
\bigr]
=
(\overline\varphi',
 \overline\varphi'')_*
\bigl[
o_{\lambda_{i+1}}(\overline\varphi)
\bigr].
\]
这里保留原书在障碍类输入中的
带横线记号
\translatornote{原书实体页 304
此前只对模 $T^{\lambda_i}$
的滤过 $\varphi$
定义 $o_{\lambda_{i+1}}$，
此处却写
$o_{\lambda_{i+1}}(\overline\varphi)$。
仅凭约化同态不足以确定下一能量层障碍；
应用此式还需相应滤过提升。
原书未交代提升记号，故保留原式；
见疑误 A-0292。}。
\end{enumerate}
\end{proposition}

命题 5.2.55 的证明与
定义--命题 4.5.16 类似，
原书从略。

现在用命题 5.2.55
完成定理 5.2.35 的滤过版本的证明。
在该定理的情形下，
对 $i$ 归纳可得下列命题。

\begin{proposition}{5.2.57}\label{prop:5-2-57}
对任意 $i=1,2,\ldots$，
存在位于
$((\mathfrak f^{(1)})^{-1},
  (\mathfrak f^{(0)})^{-1})$
之上的模 $T^{\lambda_i}$
的滤过 $A_\infty$ 双模同态
$\varphi':D'\to D$。
还存在位于
$\mathfrak F^{(i)}$
之上的模 $T^{\lambda_i}$
的滤过 $A_\infty$ 双模同态
\translatornote{原书实体页 305
此处写 “module $T^{\lambda_i}$”，
而定义 5.2.54 及本命题前一句
均写 “modulo $T^{\lambda_i}$”；
这里按同余截断译作“模”。
见疑误 A-0288。}
$\mathfrak H:D\to\mathfrak D$；
其中 $\mathfrak F^{(i)}$
是从恒等同态到
$(\mathfrak f^{(i)})^{-1}
\circ\mathfrak f^{(i)}$
的同伦，且
\[
\mathrm{Eval}_{s=0}\circ\mathfrak H
=\mathrm{id},
\qquad
\mathrm{Eval}_{s=1}\circ\mathfrak H
=\varphi'\circ\varphi.
\]
\end{proposition}

由命题 5.2.55，
证明与第 4.5.5 小节的论证相同，
原书从略。
至此，定理 5.2.35 的证明完成。
